\documentclass[aspectratio=169, xcolor=dvipsnames]{beamer}

\useoutertheme{miniframes}
\useinnertheme{circles}
\usepackage{appendixnumberbeamer}
\setbeamertemplate{footline}[frame number]

\definecolor{main_color}{HTML}{7D9D33}
\definecolor{secondary_color}{HTML}{CED38C}

\setbeamercolor{palette primary}{bg=secondary_color,fg=white}
\setbeamercolor{palette secondary}{bg=secondary_color,fg=white}
\setbeamercolor{palette tertiary}{bg=secondary_color,fg=white}
\setbeamercolor{palette quaternary}{bg=secondary_color,fg=white}
\setbeamercolor{structure}{fg=main_color}
\setbeamercolor{section in toc}{fg=secondary_color}
\setbeamercolor{subsection in head/foot}{bg=main_color,fg=white}

\usepackage{booktabs}
\usepackage{graphicx}
\usepackage{subcaption}
\usepackage{caption}

% Set the path for figures
\graphicspath{{/Users/ohouck/OneDrive - The University of Chicago/ai_weather_ag/figures/}}

\title{Tailoring Machine Learning Weather Predictions for Local Impacts}
\author{\textbf{Ozma Houck, James Franke}}
\date{}

\usepackage{bbm}

\begin{document}

\maketitle

\begin{frame}{What's Been Happening in Weather Forecasting?}
    \begin{itemize}
        \item Traditional numerical weather prediction (NWP) models are expensive and slow to run. 
        %These are physics based models which create predictions by solving the PDE's that govern the atmosphere.
        \begin{itemize}
            \item Ex: Best global weather model created by Eurpean Centre for Medium-Range Weather Forecasts (ECMWF) requires a supercomputer.
        \end{itemize}
        \pause
        \item Beginning in 2022, several machine learning (ML) based weather models have been created that have rivaled or exceeded the accuracy of traditional models.
        \item Critically, these models can be run on a personal laptop. 
    \end{itemize}
\end{frame}

\begin{frame}{Increased Accessibility of Weather Predictions Opens New Opportunities}
    \begin{itemize}
        \item Individual countries and smaller organizations can now run their own weather models instead of relying on large government agencies.
        \item Global models are more accurate in some areas than others.
        \item Different groups prioritize accuracy for different weather variables.
    \end{itemize}
        
    \pause
    \begin{itemize}
        \item \textbf{The Rub}: All global weather models (NWP or ML) are trained to optimize overall global accuracy.
        \begin{itemize}
            \item A temperature error in Greenland is treated the same as a temperature error in Chicago!
            \item The ML models might be cheap to train, but they are extremely expensive to train!
        \end{itemize}
     \item \textbf{Potential Solution}: Post-process global model outputs with a small, localized model to improve accuracy for specific regions and weather variables of interest.
    \end{itemize}
\end{frame}

\begin{frame}{How to Post-Process Weather Model Outputs? What Outputs Matter Most?}

    \begin{itemize}
        \item Just as ML tools can be used to post-process models just as they have been used to build them. 
        \begin{itemize}
            \item \textbf{Contribution}: Demonstrate that a simple low-compute neural network model can meaningfully improve the accuracy of both traditional and ML weather models. Additionally, show where these simple models perform well, and where they struggle.
            \item Additionally, show that post-processing using alternative loss functions can lead to greater improvements for specific weather variables of interest.
        \end{itemize}
    \end{itemize}

\end{frame}

\begin{frame}{Start by Manually Looking at Five Regions of Interest}
    \begin{itemize}
        \item Manually choose a few regions with a variety of climates and land terrains. 
        \item Sample from different climate zones. 
    \end{itemize}
    
    \begin{figure}
        \centering
        \includegraphics[width=0.8\linewidth]{finetuning/climate_zones_with_patches.png}
    \end{figure}
\end{frame}

\begin{frame}{See Greater Improvements for Longer Lead Times and for Wind Speed}
\begin{figure}[ht]
    \centering
    \begin{subfigure}[t]{0.48\textwidth}
        \centering
        \includegraphics[width=\linewidth]{finetuning/pangu/lead_time/multi_region/2x2/leadtime_improvement_2m_temperature_trainedwith_2m_temperature_geographic_pangu_nn_mlp.png}
        \caption{2m Temperature: Pangu (ML Model)}
    \end{subfigure}
    \hfill
    \begin{subfigure}[t]{0.48\textwidth}
        \centering
        \includegraphics[width=\linewidth]{finetuning/pangu/lead_time/multi_region/2x2/leadtime_improvement_10m_wind_speed_trainedwith_10m_wind_speed_geographic_pangu_nn_mlp.png}
        \caption{Wind Speed: Pangu (ML Model)}
    \end{subfigure}
\end{figure}
\end{frame}

\begin{frame}{Climate Zones}
\begin{figure}[ht]
    \centering
    \begin{subfigure}[t]{0.48\textwidth}
        \centering
        \includegraphics[width=\linewidth]{/Users/ohouck/Library/CloudStorage/OneDrive-TheUniversityofChicago/ai_weather_ag/figures/finetuning/pangu/lead_time/multi_region/2x2/leadtime_improvement_2m_temperature_trainedwith_2m_temperature_climate_zones_pangu_nn_mlp_bootstrap.png}
        \caption{2m Temperature: Pangu}
    \end{subfigure}
    \hfill
    \begin{subfigure}[t]{0.48\textwidth}
        \centering
        \includegraphics[width=\linewidth]{/Users/ohouck/Library/CloudStorage/OneDrive-TheUniversityofChicago/ai_weather_ag/figures/finetuning/pangu/lead_time/multi_region/2x2/leadtime_improvement_10m_wind_speed_trainedwith_10m_wind_speed_climate_zones_pangu_nn_mlp_bootstrap.png}
        \caption{Wind Speed: Pangu}
    \end{subfigure}
\end{figure}
\end{frame}

\end{document}


